\section{Démonstrations techniques}
\label{app:demonstrations_techniques}

\subsection{NP-difficulté de la sélection des gardiens}
\label{app:preuve_np_difficulte}

\begin{proof}[Démonstration de la Proposition~\ref{prop:np_difficulte_gardiens}]
Considérons une instance de \textsc{Subset Sum}, problème faiblement
NP-complet \cite{garey-johnson-1979}, composée de $m$ entiers strictement
positifs $u_1,\ldots,u_m$ et d'une cible $B>0$. Le cas \(m=1\) étant
trivial, supposons \(m\geqslant2\). Posons
$U:=\sum_{i=1}^m u_i$ (si $B>U$, la réponse est immédiatement négative)
et supposons donc $B\leqslant U$.

On fixe une fenêtre réduite à une heure, \(H=\{h\}\), puis on construit
$N:=\{1,\ldots,m\}$ et la coalition $S:=N$ en posant
\[
q_i:=Uu_i,\qquad
F_i:=u_i,\qquad
\beta_i:=u_i,\qquad
d_i^h:=Bu_i.
\]
Les paramètres primitifs produits sont entiers, leur taille d'encodage est
polynomiale, et $\gamma_i=\beta_i/q_i=1/U$ pour tout $i\in N$. De plus,
$d_i^h\leqslant q_i$ et
\[
d^h(S)=\sum_{i\in S}d_i^h=BU.
\]
Cette construction s'effectue en temps polynomial dans la taille de
l'instance initiale.

Pour tout $G\subseteq N$, la configuration est faisable si et seulement si
\[
q(G)=U\sum_{i\in G}u_i\geqslant d^h(S)=BU,
\]
soit si et seulement si $\sum_{i\in G}u_i\geqslant B$. Comme tous les
coûts variables unitaires valent $1/U$, toute répartition admissible a
pour coût variable
\[
\sum_{i\in G}\gamma_i t_i^h
=\frac{1}{U}d^h(S)=B.
\]
Par conséquent,
\[
C_H^*(S)
=B+\min\left\{
\sum_{i\in G}u_i\ \middle|\
G\subseteq N,\ \sum_{i\in G}u_i\geqslant B
\right\}.
\]

Puisque tout ensemble du minimum a une somme au moins égale à $B$,
\[
C_H^*(S)\leqslant2B
\quad\Longleftrightarrow\quad
\min_{\substack{G\subseteq N\\ \sum_{i\in G}u_i\geqslant B}}
\sum_{i\in G}u_i=B
\quad\Longleftrightarrow\quad
\exists\,G\subseteq N,\ \sum_{i\in G}u_i=B.
\]
Un algorithme polynomial pour la sélection optimale des gardiens
permettrait donc de décider \textsc{Subset Sum} en temps polynomial. Le
problème de sélection optimale des gardiens est ainsi NP-difficile.
\end{proof}

\subsection{Complexité de l'énumération exacte}
\label{app:preuve_complexite_enumeration}

\begin{proof}[Démonstration de la Proposition~\ref{prop:complexite_enumeration}]
Après avoir ordonné les membres de \(S\) par coût variable unitaire, une
coalition de taille \(s\) admet \(2^s-1\) ensembles de gardiens. Pour chacun,
la faisabilité et la répartition gloutonne sont évaluées aux \(|H|\) heures en
\(O(|H|s)\). En sommant sur les coalitions, on obtient
\[
|H|\sum_{s=1}^{n}\binom ns s(2^s-1)
=|H|\left(2n\,3^{n-1}-n\,2^{n-1}\right)
=O(|H|n\,3^n).
\]
\end{proof}

\subsection{Caractérisation du cas homogène}
\label{app:preuve_cas_homogene}

\begin{proof}[Démonstration de la Proposition~\ref{prop:coeur_cas_homogene}
et du Corollaire~\ref{cor:partages_cas_homogene}]
Une coalition de taille \(s\) requiert au moins \(k_{H,s}\) gardiens pour
rester faisable à toutes les heures, et tout ensemble de \(k_{H,s}\) gardiens
est faisable. Comme le coût variable total sur la fenêtre vaut toujours
\(\gamma sD_H\), une solution optimale active exactement \(k_{H,s}\)
équipements. Les expressions de \(C_H^*(S)\) et de \(v_H(S)\) en découlent.

Si \(z\in\mathcal C(N,v_H)\), la somme des contraintes du coeur sur toutes les
coalitions de taille \(s\) donne
\[
\binom{n-1}{s-1}v_{H,n}
=
\sum_{\substack{S\in\mathcal S\\|S|=s}}z(S)
\geqslant
\binom{n}{s}v_{H,s},
\]
soit \(v_{H,s}/s\leqslant v_{H,n}/n\). Réciproquement, sous ces inégalités,
l'allocation \(z_i=v_{H,n}/n\) satisfait toutes les contraintes du coeur.
La symétrie et l'efficacité donnent les expressions de Shapley et du
nucléole. Enfin, pour toute allocation efficace, l'excès moyen des coalitions
de taille \(s\) vaut \(v_{H,s}-sv_{H,n}/n\). L'allocation égalitaire rend tous
ces excès égaux à leur moyenne, ce qui établit la formule du moindre coeur.
\end{proof}

\subsection{Certificat leave-one-out}
\label{app:preuve_leave_one_out}

\begin{proof}[Démonstration de la Proposition~\ref{prop:certificat_leave_one_out}
et du Corollaire~\ref{cor:borne_leave_one_out}]
Attribuons le poids \(1/(n-1)\) à chaque coalition
\(N\setminus\{i\}\), et un poids nul aux autres coalitions. Cette famille est
équilibrée. En utilisant la définition de \(v_H\), on obtient
\begin{align*}
    \frac{1}{n-1}\sum_{i\in N}v_H(N\setminus\{i\})-v_H(N)
    &=
    C_H^*(N)
    -
    \frac{1}{n-1}
    \sum_{i\in N}C_H^*(N\setminus\{i\})\\
    &=
    \Delta_{\mathrm{LOO},H}(N).
\end{align*}
Ainsi, \(\Delta_{\mathrm{LOO},H}(N)>0\) viole la condition de
Bondareva--Shapley.

Considérons maintenant toute solution admissible \((z,\varepsilon)\) du
programme \eqref{eq:lp_moindre_coeur}. En sommant ses contraintes pour les
\(n\) coalitions \(N\setminus\{i\}\), il vient
\[
    (n-1)v_H(N)
    =
    \sum_{i\in N}z(N\setminus\{i\})
    \geqslant
    \sum_{i\in N}v_H(N\setminus\{i\})-n\varepsilon.
\]
Par conséquent,
\[
    \varepsilon
    \geqslant
    \frac{1}{n}
    \left[
        \sum_{i\in N}v_H(N\setminus\{i\})-(n-1)v_H(N)
    \right]
    =
    \frac{n-1}{n}\Delta_{\mathrm{LOO},H}(N).
\]
Cette borne valant pour toute solution admissible, elle vaut en particulier
pour \(\varepsilon^*\).
\end{proof}

\subsection{Capacité individuelle suffisante sur une fenêtre}
\label{app:preuve_capacite_coeur}
\label{app:preuve_capacite_fenetre}

\begin{proof}[Démonstration du Théorème~\ref{thm:capacite_coeur_non_vide}]
Soit \(Q\in\mathcal S\). Sous l'hypothèse de capacité,
tout \(j\in Q\) peut servir seul \(Q\) à chaque heure. Réciproquement,
considérons un ensemble persistant de gardiens \(G\subseteq Q\) et choisissons
\(j\in\operatorname*{arg\,min}_{i\in G}\gamma_i\). Pour toute famille de
répartitions faisables \((\mathbf t^h)_{h\in H}\),
\begin{align*}
    \sum_{h\in H}\sum_{i\in G}\left(F_i+\gamma_i t_i^h\right)
    &\geqslant
    |H|F_j+\gamma_j\sum_{h\in H}\sum_{i\in G}t_i^h\\
    &=|H|F_j+\gamma_jD(Q).
\end{align*}
La configuration persistante réduite au gardien \(j\) est faisable. Il en
résulte
\[
    C_H^*(Q)
    =
    \min_{j\in Q}
    \left\{|H|F_j+\gamma_jD(Q)\right\}.
\]
Posons \(p_H:=C_H^*(N)/D(N)\) et \(y_{i,H}:=p_HD_i\). Pour tout
\(Q\in\mathcal S\) et tout \(j\in Q\),
\[
    p_H
    \leqslant
    \gamma_j+\frac{|H|F_j}{D(N)}
    \leqslant
    \gamma_j+\frac{|H|F_j}{D(Q)}.
\]
Ainsi, \(\sum_{i\in Q}y_{i,H}\leqslant C_H^*(Q)\) et
\(\sum_{i\in N}y_{i,H}=C_H^*(N)\). Pour \(Q=\{i\}\), la première inégalité
donne \(z_i=C_H^*(\{i\})-y_{i,H}\geqslant0\), tandis que
\(\sum_{i\in N}z_i=v_H(N)\). La
Proposition~\ref{prop:translation_coeur} permet donc de conclure que
\(z\in\mathcal C(N,v_H)\).
\end{proof}

\subsection{Instabilité de Shapley en très faible trafic}
\label{app:preuve_shapley_instable}

\begin{proof}[Démonstration de la Proposition~\ref{prop:shapley_hors_coeur_tres_faible_trafic}]
On a \(d^h(N)=4<10=q_i\) pour tout \(i\in N\) et tout \(h\in H\) :
chaque équipement peut donc servir seul la demande de la grande coalition à
chaque heure. Les coûts autonomes sur la fenêtre sont
\[
    \bigl(C_H^*(\{1\}),C_H^*(\{2\}),C_H^*(\{3\}),C_H^*(\{4\})\bigr)
    =|H|(2,2,5,5).
\]
Comme \(d^h(S)=|S|\) et que tous les coûts variables unitaires valent un,
\eqref{eq:cout_tres_faible_trafic} donne, pour tout \(S\in\mathcal S\),
\[
    C_H^*(S)=|H|\left(|S|+\min_{i\in S}F_i\right),
    \qquad
    v_H(S)=|H|\left(\sum_{i\in S}F_i-\min_{i\in S}F_i\right).
\]
En particulier,
\[
    v_H(N)=9|H|,
    \qquad
    v_H(\{1,3,4\})=8|H|.
\]

Les opérateurs 1 et 2 sont symétriques, de même que les opérateurs 3 et 4.
Pour l'opérateur 3, la contribution marginale est nulle lorsqu'il arrive en
premier et vaut \(4|H|\) dès qu'au moins un opérateur le précède. Il vient donc
\[
    \phi_3(v_H)=\phi_4(v_H)=3|H|.
\]
Pour l'opérateur 1, trois événements disjoints suffisent. Il arrive en premier
avec probabilité \(1/4\), auquel cas sa contribution est nulle ; l'opérateur 2
le précède avec probabilité \(1/2\), auquel cas sa contribution vaut \(|H|\) ;
enfin, il précède l'opérateur 2 sans arriver en premier avec probabilité
\(1/2-1/4=1/4\), auquel cas une coalition formée uniquement d'opérateurs à
coût fixe élevé le précède et sa contribution vaut \(4|H|\). Par conséquent,
\[
    \phi_1(v_H)=\phi_2(v_H)
    =|H|\left(\frac14\,0+\frac12\,1+\frac14\,4\right)
    =\frac32|H|.
\]
L'efficacité est vérifiée puisque
\(|H|\bigl(2(3/2)+2(3)\bigr)=9|H|=v_H(N)\). Pour
\(Q=\{1,3,4\}\),
\[
    \sum_{i\in Q}\phi_i(v_H)
    =\frac{15}{2}|H|
    <8|H|
    =v_H(Q).
\]
La contrainte du cœur associée à \(Q\) est donc violée et
\(\phi(v_H)\notin\mathcal C(N,v_H)\).

Le cœur n'est pourtant pas vide. L'allocation du
Théorème~\ref{thm:capacite_coeur_non_vide} vaut ici
\[
    z^{\mathrm{core}}
    =|H|\left(\frac34,\frac34,\frac{15}{4},\frac{15}{4}\right),
    \qquad
    z^{\mathrm{core}}(N)=9|H|=v_H(N).
\]
Pour vérifier toutes ses contraintes, posons
\(a:=|S\cap\{1,2\}|\) et \(b:=|S\cap\{3,4\}|\). Si \(a\geqslant1\), alors
\[
    z^{\mathrm{core}}(S)-v_H(S)
    =|H|\left[\frac{3a+15b}{4}-(a+4b-1)\right]
    =|H|\frac{4-|S|}{4}
    \geqslant0.
\]
Si \(a=0\) et \(b\geqslant1\), alors
\[
    z^{\mathrm{core}}(S)-v_H(S)
    =|H|\left[\frac{15b}{4}-4(b-1)\right]
    =|H|\left(4-\frac b4\right)
    >0.
\]
Ainsi \(z^{\mathrm{core}}\in\mathcal C(N,v_H)\), ce qui distingue bien la
non-appartenance de Shapley de la non-vacuité du cœur.
\end{proof}

\subsection{Stabilité de Shapley sous homogénéité des coûts}
\label{app:preuve_shapley_homogene}

\begin{proof}[Démonstration du Corollaire~\ref{cor:shapley_coeur_tres_faible_trafic_homogene}]
Par \eqref{eq:cout_tres_faible_trafic},
\(C_H^*(S)=|H|F+\gamma D(S)\) pour tout \(S\in\mathcal S\). Comme
\(C_H^*(\{i\})=|H|F+\gamma D_i\),
\[
    v_H(S)
    =\sum_{i\in S}(|H|F+\gamma D_i)-|H|F-\gamma D(S)
    =|H|F(|S|-1).
\]
Le jeu ne dépend que du cardinal des coalitions. Par symétrie et efficacité,
si \(n=|N|\),
\[
    \phi_i(v_H)=\frac{v_H(N)}{n}=\frac{|H|F(n-1)}{n},
    \qquad \forall i\in N.
\]
Pour tout \(Q\in\mathcal S\) de cardinal \(k\),
\[
    \sum_{i\in Q}\phi_i(v_H)-v_H(Q)
    =\frac{k|H|F(n-1)}{n}-|H|F(k-1)
    =\frac{|H|F(n-k)}{n}
    \geqslant0.
\]
L'efficacité de la valeur de Shapley et les contraintes précédentes donnent
\(\phi(v_H)\in\mathcal C(N,v_H)\).
\end{proof}

\subsection{Équilibre budgétaire des transferts}
\label{app:preuve_equilibre_budgetaire}

\begin{proof}[Démonstration de la Proposition~\ref{prop:equilibre_budgetaire}]
\begin{align*}
\sum_{i\in N}\tau_i
&=
\sum_{i\in N}C_{i,H}^*(N)
-
\sum_{i\in N}C_H^*(\{i\})
+
\sum_{i\in N}z_i\\
&=
C_H^*(N)-\sum_{i\in N}C_H^*(\{i\})+v_H(N)
=
0.
\end{align*}
\end{proof}

\subsection{Équivalence entre économies attribuées et coûts nets}
\label{app:preuve_translation_coeur}

\begin{proof}[Démonstration de la Proposition~\ref{prop:translation_coeur}]
L'efficacité de \(z\) donne
\[
\sum_{i\in N}y_i
=
\sum_{i\in N}C_H^*(\{i\})-v_H(N)
=
C_H^*(N).
\]
Pour tout \(Q\in\mathcal S\),
\begin{align*}
\sum_{i\in Q}y_i\leqslant C_H^*(Q)
&\Longleftrightarrow
\sum_{i\in Q}C_H^*(\{i\})-z(Q)\leqslant C_H^*(Q)\\
&\Longleftrightarrow
z(Q)\geqslant
\sum_{i\in Q}C_H^*(\{i\})-C_H^*(Q)\\
&\Longleftrightarrow
z(Q)\geqslant v_H(Q).
\end{align*}
Les contraintes de stabilité sont donc équivalentes.
\end{proof}
